% Worked Examples: Concept 3.7.5 - Implementation
\documentclass[11pt,a4paper]{article}
\usepackage{worked-examples-template}

\title{\textbf{Worked Examples}\\[0.5em]
\large Concept 3.7.5: Implementation\\[0.3em]
\normalsize \coursesubtitle{} -- \coursetitle}
\author{Assoc.\ Prof.\ William Robertson \and Dr Sean McGowan}
\date{\today}

\begin{document}

\maketitle
\tableofcontents
\newpage

\section{Concept 3.7.5: Implementation}

\subsection{Example 1: Discrete-Time PID Implementation}

\paragraph{Problem:} A continuous PID controller:
\begin{align}
u(t) = K_p e(t) + K_i \int_0^t e(\tau)d\tau + K_d \frac{de(t)}{dt}
\end{align}

with $K_p = 2$, $K_i = 0.5$, $K_d = 1$ needs to be implemented digitally with sampling period $T_s = 0.1$ seconds.

(a) Derive the discrete-time implementation using:
\begin{itemize}
\item Forward Euler for the integral
\item Backward difference for the derivative
\end{itemize}

(b) Write the resulting difference equation.

\paragraph{Solution:}

\textbf{Part (a): Discretization}

Let $e_k = e(kT_s)$ denote the error at time step $k$.

\textbf{Proportional term:} No approximation needed
\begin{align}
u_P[k] = K_p e_k = 2e_k
\end{align}

\textbf{Integral term:} Using forward Euler (rectangular rule)
\begin{align}
\int_0^{kT_s} e(\tau)d\tau \approx T_s \sum_{j=0}^{k} e_j
\end{align}

Define the accumulated sum:
\begin{align}
I_k = I_{k-1} + T_s e_k
\end{align}

Then:
\begin{align}
u_I[k] = K_i I_k = 0.5 I_k
\end{align}

\textbf{Derivative term:} Using backward difference
\begin{align}
\frac{de}{dt}\bigg|_{t=kT_s} \approx \frac{e_k - e_{k-1}}{T_s}
\end{align}

Therefore:
\begin{align}
u_D[k] = K_d \frac{e_k - e_{k-1}}{T_s} = 1 \cdot \frac{e_k - e_{k-1}}{0.1} = 10(e_k - e_{k-1})
\end{align}

\textbf{Part (b): Complete Difference Equation}

The total control signal is:
\begin{align}
u_k = u_P[k] + u_I[k] + u_D[k]
\end{align}

Substituting:
\begin{align}
u_k = 2e_k + 0.5I_k + 10(e_k - e_{k-1})
\end{align}

With the integral update:
\begin{align}
I_k = I_{k-1} + 0.1e_k
\end{align}

Combining:
\begin{align}
u_k = 2e_k + 0.5(I_{k-1} + 0.1e_k) + 10e_k - 10e_{k-1}
\end{align}

Simplifying:
\begin{align}
u_k = 0.5I_{k-1} + (2 + 0.05 + 10)e_k - 10e_{k-1}
\end{align}

\begin{align}
u_k = 0.5I_{k-1} + 12.05e_k - 10e_{k-1}
\end{align}

\textbf{Algorithm:}

\begin{verbatim}
Initialize: I = 0, e_prev = 0

At each time step k:
1. Read measurement y_k
2. Compute error: e_k = r_k - y_k
3. Compute control:
   u_k = 0.5*I + 12.05*e_k - 10*e_prev
4. Update integral: I = I + 0.1*e_k
5. Update previous error: e_prev = e_k
6. Apply control signal: u_k (with saturation if needed)
\end{verbatim}

\textbf{Alternative Velocity Form:}

For better numerical properties and easier anti-windup:
\begin{align}
\Delta u_k = u_k - u_{k-1} = q_0 e_k + q_1 e_{k-1} + q_2 e_{k-2}
\end{align}

where:
\begin{align}
q_0 &= K_p + K_i T_s + \frac{K_d}{T_s} = 2 + 0.05 + 10 = 12.05 \\
q_1 &= -K_p - 2\frac{K_d}{T_s} = -2 - 20 = -22 \\
q_2 &= \frac{K_d}{T_s} = 10
\end{align}

This form computes the change in control signal:
\begin{align}
u_k = u_{k-1} + 12.05e_k - 22e_{k-1} + 10e_{k-2}
\end{align}

\textbf{Practical Considerations:}
\begin{itemize}
\item Add low-pass filter on derivative to reduce noise sensitivity
\item Implement anti-windup for integral term
\item Use appropriate data types (floating-point for better precision)
\item Consider derivative on measurement rather than error to avoid derivative kick
\end{itemize}
\end{document}
